../cictikz/src/cictikz/data/tex/cictikz_preamble.tex

% First-order Gm-C section.  There is only one node: G_m1's output, C_A,
% C_X and G_m2's output all meet there, and the rail carries it on to V_O.
% G_m2 is drawn mirrored -- it senses V_O on its tall right-hand edge and
% drives the node from its short left-hand edge -- so it loads the node with
% a conductance G_m2.
%
%   s C_A V_o = G_m1 V_i + s C_X (V_i - V_o) - G_m2 V_o
%   V_o/V_i = [s C_X/(C_A+C_X) + G_m1/(C_A+C_X)] / [s + G_m2/(C_A+C_X)]
%
% which is the H(s) the lecture prints.  The i_0..i_3 blanks under the
% drawing are the exercise the original poses; they are kept.

\begin{circuitikz}[american, thick, transform shape]

  ../cictikz/src/cictikz/data/tex/cictikz_lib.tex
  ../cictikz/src/cictikz/data/tex/gmc_lib.tex

  \def\ynode{2.8}   % the one node, and the V_O rail
  \def\xnode{4.35}

  % ---- G_m1 ----
  \gmcBody{1.1}{\ynode}{$G_{m1}$}
  \draw (0.36,3.6) -- (1.1,3.6);
  \gmcPlus{1.45}{3.6}
  \draw (0.36,2.0) -- (1.1,2.0);
  \gmcMinus{1.45}{2.0}
  \draw (0.36,2.0) -- (0.36,1.45);
  \draw (0.36,1.45) \vground;
  \node[anchor=east] at (0.3,3.6) {$V_{i}(s)$};

  % ---- C_X feeds the node directly from the input ----
  \draw (0.72,3.6) -- (0.72,5.0);
  \gmcHcap{0.72}{5.0}{2.83}
  \node[anchor=south] at (1.78,5.42) {$C_{X}$};
  \draw (2.83,5.0) -- (\xnode,\ynode);

  % ---- G_m1 into the node ----
  \draw (2.6,\ynode) -- (\xnode,\ynode);
  \fill (\xnode,\ynode) circle (0.075);

  % ---- C_A on the node ----
  \gmcVcap{\xnode}{\ynode}{0.75}
  \draw (\xnode,0.75) \vground;
  \node[anchor=west] at (4.75,1.78) {$C_{A}$};

  % ---- G_m2, mirrored: senses V_O, drives the node ----
  \gmcBodyL{8.77}{5.0}{$G_{m2}$}
  \gmcPlus{8.45}{5.8}
  \gmcMinus{8.45}{4.2}
  \draw (8.77,5.8) -- (9.5,5.8) -- (9.5,5.45);
  \draw (9.5,5.45) \vground;
  \draw (8.77,4.2) -- (9.5,4.2) -- (9.5,\ynode);
  \draw (7.27,5.0) -- (\xnode,\ynode);

  % ---- the output rail ----
  \draw (\xnode,\ynode) -- (11.2,\ynode);
  \fill (9.5,\ynode) circle (0.075);
  \node[anchor=west] at (11.28,\ynode) {$V_{O}(s)$};
  \node[anchor=north west] at (11.0,2.6) {$+$};
  \node[anchor=north west] at (11.0,1.75) {$-$};
  \draw (11.15,1.25) -- (11.15,1.02);
  \draw (11.15,1.02) \vground;

  % ---- the four branch currents the exercise asks for ----
  \draw[->] (3.11,2.6) -- (3.84,2.6);
  \node[anchor=north] at (3.45,2.52) {$i_{0}$};
  \draw[->] (3.94,3.73) -- (4.39,3.27);
  \node[anchor=west] at (4.5,3.55) {$i_{1}$};
  \draw[->] (5.61,3.48) -- (5.15,3.12);
  \node[anchor=west] at (5.72,3.32) {$i_{2}$};
  \draw[->] (5.0,2.6) -- (5.61,2.6);
  \node[anchor=north] at (5.3,2.52) {$i_{3}$};

  % ---- the blanks the original leaves for the student ----
  \node[anchor=west] at (1.2,-0.6) {$i_{0} =$};
  \node[anchor=west] at (1.2,-1.3) {$i_{1} =$};
  \node[anchor=west] at (1.2,-2.0) {$i_{2} =$};
  \node[anchor=west] at (1.2,-2.7) {$i_{3} =$};

\end{circuitikz}
\end{document}
